\section{Transferring a strong Pell root into a transport gap}
\label{sec:root-transfer}

This section gives a polynomial in twelve strictly positive integer
unknowns of exact degree $84$, with a literal $103$-operation arithmetic
circuit.  Its predecessor in package 61 has degree $118$ and $104$
operations.  A thirteen-positive-variable companion has degree $30$ and
$94$ operations, against package 61's degree $34$ and $93$ operations.
These comparisons concern the actual circuits, count all live binary
additions, subtractions and multiplications, and treat the complete fixed
coefficient numerals as leaves.

The new twelve-variable construction removes the normalized strong Pell
root through the first transport carrier.  A positive upper approximation
to that root makes both recovered root signs harmless.  Strengthening an
existing transport multiplier makes the inverse coordinate positive.
Thus the construction does not require a new index congruence for the
auxiliary Pell ordinate.  The smaller thirteen-variable companion combines
package 61's transport carrier with the adaptive modulus of package 65.

\subsection{Fixed coefficients and the precise inherited boundary}

Use the whole compact coefficient tuple of package 44:
\[
 \beta=2^{590069}-1,\qquad K_0,\qquad M_C,\qquad M_F.
\]
The four integers are delivered as lossless payloads in
\texttt{coefficients/}; $M_F$ is already the shifted source mask.
They are constants, not existential parameters.  The fixed numerical
binding and the unchanged source proof are recorded separately from the
new reductions.

We use package 61's normalized positive-source interface.  With positive
coordinates $J,F,\alpha,z,f,h,i,T_{\rm old},s,w,\tau,\eta,\zeta,y,
\gamma_{\rm old}$, define
\begin{align*}
 q&=\beta J+1,& C&=q-F-\alpha,\\
 X&=wq,&Y&=sq^3,\\
 E&=XY,& k&=\eta+\zeta,\\
 c&=kY+\eta,&a&=E+Y,\\
 H&=4a+3,&\Delta&=a^2+H=(a+2)^2-1,\\
 \ell&=EkY,&I&=ic^2.
\end{align*}
\[
 R=\big((q-1)(q-F)+\alpha\big)(q^2-1)+(M_C+qM_F)J.
\]
The six source equations are
\begin{align}
 \tau^2-\ell(\ell+k)&=1,&
 D^2-\Delta c^2&=1,\qquad D=X+ac+H\gamma_{\rm old},\label{eq:rt-source-a}\\
 f^2-\Delta I^2&=1,&
 \Delta(f^2-1)(V^2-y^2)+y^2&=1,\label{eq:rt-source-b}\\
 k-hE-R&=1,&
 (K_0+w)C+(q-F)-z(q-1)&=1,\label{eq:rt-source-c}
\end{align}
where $V=c(T_{\rm old}f-1)-Rf^2$.
The inherited compiler and K215 calibration identify existence of a
positive solution with a finite syntactic ZFC refutation.  This source
interface, rather than a finite experiment, is the semantic premise used
below.

Put $B=\beta+1$.  We use the source's canonical completeness promises
\[
 q=B^N,\quad J=1+B+\cdots+B^{N-1},\quad
 N=581197^e>1,\quad R\equiv3\pmod4,\quad X=2^R.
\]
The first Pell index is $n_0=(R+1)/2\ge3$.  A normalized strong solution
can be chosen together with a positive auxiliary solution of the form
\[
 V=f^2(cT-R)-c,\qquad y=V+u,\qquad T,u>0.
\]
This is the normalized source and canonical construction of package 61's report, Sections 2 and 6, together with its positive auxiliary reconstruction lemma (Section 3, source label \texttt{lem:auxcomplete});
no stronger congruence for $y$ is required here.

For arbitrary positive outer coordinates, once $C>0$ has been proved,
the authentic masks give
\begin{equation}
 q\ge16,\qquad 3q<R<q^4-2.\label{eq:rt-Rbound}
\end{equation}
These bounds do not presuppose $q\mid X$.

\subsection{A weighted transport carrier and the adaptive main carrier}

Initially supply $X,Y$ directly.  Let $j,F,\alpha,f,i,T,X,Y,r,\eta,
\zeta,u,\gamma$ be positive, put $J=3j+1$, and define the registers above.
Let $h_*>0$ be an integer polynomial register, and set
\begin{align}
 W&=1+(q-1)h_*,&M&=q^3(q-1),\nonumber\\
 B_t&=(qK_0+X)C-qF,&
 Z&=M(\ell+k+r)+(q-1)Y-q^3WB_t,\label{eq:rt-transport}\\
 A_0&=Z^2-M^2\bigl(\ell(\ell+k)+1\bigr).\nonumber
\end{align}
The adaptive main carrier is
\begin{align}
 S&=4qY,& U&=qS=4q^2Y,& r_0&=k-R-1,\nonumber\\
 D_*&=U(X+ac)+H(U\gamma+SX-r_0),\label{eq:rt-main}\\
 A_1&=D_*^2-U^2(1+\Delta c^2).\nonumber
\end{align}
Keep the other two residuals
\begin{align}
 A_s&=f^2-\Delta I^2-1,\nonumber\\
 V&=f^2(cT-R)-c,\nonumber\\
 A_a&=(V+u)^2-\Delta(f^2-1)u(2V+u)-1.\label{eq:rt-aux}
\end{align}

\paragraph{Weighted-carrier lemma.}
Every positive simultaneous zero of $A_0,A_1,A_a,A_s$, with any positive
integer $h_*$, reconstructs a positive solution of
\eqref{eq:rt-source-a}--\eqref{eq:rt-source-c}.  Conversely, canonical
source witnesses lift whenever the chosen multiplier satisfies $h_*\ge k$.

\begin{proof}
If a rational number has integral square, it is an integer.  Apply this
to $Z/M$: $A_0=0$ supplies an integer $\tau$ with
$\tau^2=\ell(\ell+k)+1$.  Since $M\mid Z$, reduction modulo $q^3$
gives $q^3\mid Y$.  Reduction modulo $q-1$ gives $(q-1)\mid B_t$,
because $W\equiv q^3\equiv1\pmod{q-1}$.  Define the proof-only integers
$s=Y/q^3>0$ and $t=B_t/(q-1)$.  Dividing the carrier gives
\[
 \tau=\ell+k+r+s-Wt.
\]
As $|\tau|<\ell+k$, we obtain $Wt>r+s>0$, hence $t>0$ and $B_t>0$.
If $C\le0$, then $B_t=(qK_0+X)C-qF<0$, a contradiction.  Thus $C>0$
and \eqref{eq:rt-Rbound} is now available.

Write $A=a+2$ and $P_0=2XY^2+1$.  The main numerator is positive before
either $q\mid X$ or rank recovery: since $r_0<k$, $c>kY$ and
$UaY>H$, we have $D_*>Uac-Hk>0$.  Thus $A_1=0$ gives a positive integer
$D=D_*/U$ with $D^2-\Delta c^2=1$.
Pell classification supplies
\[
 D=\chi_A(p),\quad c=\psi_A(p),\qquad
 |\tau|=\chi_{P_0}(n_0),\quad k=2\psi_{P_0}(n_0).
\]
Here $p\ge2$ and $n_0\ge2$.  Indeed $n_0=1$ would give $k=2$ and
$c=2Y+1<2A=\psi_A(2)$, which is impossible.  Consequently
\[
 k>8XY^2,\qquad c>8XY^3\ge8q^9>2R,\qquad n_0<p,\qquad c>2p.
\]
The normalized strong residual implies $f=\chi_A(m)$,
$I=\psi_A(m)$, $c\mid m$, and $m>2p$; this is the normalized-rank
lemma of package 61.  Put $S_*=\Delta I$.  The auxiliary residual gives
\[
 S_*^2V^2-(S_*^2-1)y^2=1,\qquad y=|V+u|>0,
\]
with $V\equiv-c\pmod f$ and $V\equiv-R\pmod c$.
The signed auxiliary-index lemma, with $0<p,R<c/2$, yields $p=R$.
The rank lemmas used here do not require the deleted scale or index
quotient, nor the sign of $V$.

It remains to extract the adaptive congruence without discarding a factor
of three.  Since $590069$ is odd, $\beta\equiv1\pmod3$ and
$J=3j+1$ imply $q\equiv2\pmod3$.  Because $q^3\mid Y$,
$\gcd(q,H)=1$.  Moreover $H=4Y(X+1)+3$ is odd and
$\gcd(H,Y)=\gcd(3,Y)$.  Hence
\[
 g:=\gcd(H,U)=\gcd(H,S)\in\{1,3\}.
\]
Integrality of $D_*/U$ gives $U/g\mid SX-r_0$, hence $S/g\mid r_0$.
The first Pell recurrence gives
\[
 k\equiv2n_0\pmod{4XY^2}.
\]
Since $S\mid4XY^2$ and $S/g\ge4q^4/3>R+1$,
the inequalities $1\le n_0<p=R$ force
$2n_0-R-1=0$.  Thus $r_0=k-2n_0$ is a positive multiple of $4XY^2$.
Also $U\mid4XY^2$, since $q^2\mid Y$; therefore $U\mid r_0$.
It follows from $U/g\mid SX$ that $q\mid gX$.
Because $g$ divides $3$ and $3\nmid q$, this gives $q\mid X$.

Now $w=X/q>0$, $h=r_0/(XY)>0$, and
$\gamma_{\rm old}=\gamma+w-r_0/U$ are integers.  Since
$D>(a+1)c$ and $c>X$, the identity
$H\gamma_{\rm old}=D-ac-X$ proves $\gamma_{\rm old}>0$.
Since $B_t=q((K_0+w)C-F)>0$ and $(q-1)\mid B_t$, the integer
$z_{\mathrm{new}}=((K_0+w)C-F)/(q-1)$ is positive; set $z=z_{\mathrm{new}}+1$.
Finally replace $\tau$ by $|\tau|$, put $T_{\rm old}=fT$, and retain
$y=|V+u|$.  These are the required source coordinates.

For completeness, use canonical source witnesses.  Their odd repunit
exponent gives $J\equiv1\pmod3$, so $j=(J-1)/3>0$ is integral.
The source transport coordinate has $z\ge2$, hence
$t=q(z-1)\ge q$.  Put
\[
 r=Wt+\tau-\ell-k-s.
\]
Since $\tau>\ell$, $s<k$, and $h_*\ge k$, this gives
$r>qW-2k>0$.  It makes $A_0=0$.
The canonical first index gives $r_0$ a positive multiple of $4XY^2$;
take $\gamma=\gamma_{\rm old}-w+r_0/U$.
It is positive because
$r_0/U\ge XY/q^2=wY/q>w$.  Its substitution makes $A_1=0$.
The positive auxiliary reconstruction supplies $T,u$, and the two
remaining residuals vanish.
\end{proof}

Taking $h_*=k$ defines the thirteen-variable companion
\[
 P_{13}=A_0^2+A_1^2+A_a^2+A_s^2.
\]
Its residual degrees are respectively $15,13,13,14$, so its exact total
degree is $30$.  The factor $4$ in $S=4qY$ is used in the proof:
removing it without another argument loses the bound $S/g>R+1$.

\subsection{The twelve-variable polynomial}

Supply exactly
\[
 j,F,\alpha,i,T,X,Y,n,\eta,\zeta,u,\gamma\in\mathbb Z_{>0}.
\]
Use the same common registers, but now define
\begin{align}
 D_f&=1+\Delta I^2,& L&=(a+2)I,& h_*&=ikc,\nonumber\\
 W&=1+(q-1)ikc,&
 Z_0&=M(\ell+k+n)+(q-1)Y-q^3WB_t,\label{eq:rt-newregs}\\
 \overline Z&=Z_0+ML,& D_\tau&=\ell(\ell+k)+1.\nonumber
\end{align}
The first residual is the ordinary polynomial
\begin{equation}
 N=\bigl(\overline Z^2-M^2(D_f+D_\tau)\bigr)^2
                  -4M^4D_fD_\tau.\label{eq:rt-N}
\end{equation}
Retain $A_1$ from \eqref{eq:rt-main} and replace the auxiliary registers by
\begin{equation}
 V=D_f(cT-R)-c,\qquad Q=\Delta(D_f-1),\qquad
 A_a=(V+u)^2-Qu(2V+u)-1.\label{eq:rt-newaux}
\end{equation}
The recommended equation is
\begin{equation}
 \boxed{P_{12}=N^2+A_1^2+A_a^2=0.}\label{eq:rt-P12}
\end{equation}
Every named register is a polynomial abbreviation.  There are twelve,
not twelve plus the number of registers, quantified unknowns.

\paragraph{Root-transfer theorem.}
The polynomial \eqref{eq:rt-P12} has a strictly positive zero if and only
if the inherited source has a strictly positive zero.  Consequently its
lack of positive zeros is equivalent in ZFC to $\operatorname{Con}(\mathrm{ZFC})$,
subject exactly to the inherited compiler/calibration interface above.

\begin{proof}
The inequalities needed to recover roots hold on every positive input,
before any residual is assumed zero.  Since $\ell<ac$ and $k<c$,
\[
 \ell(\ell+k)<a(a+1)c^2<\Delta i^2c^4,
\]
and therefore $D_f>D_\tau>0$.  Also
\[
 L^2-D_f=I^2-1>0,
\]
because $c\ge3$ and $I=ic^2>1$.

At a zero of $N$, $\overline Z=0$ is impossible: it would imply
$M^4(D_f-D_\tau)^2=0$.  The identity
\begin{align*}
 N={}&\bigl(\overline Z^2+M^2(D_f-D_\tau)\bigr)^2
                     -4M^2\overline Z^2D_f
\end{align*}
then shows that the rational number
\[
 f_*=\frac{\overline Z^2+M^2(D_f-D_\tau)}{2M\overline Z}
\]
has square $D_f$.  It is a nonzero integer by rational-square
integrality.  Furthermore
\[
 \left(\frac{\overline Z-Mf_*}{M}\right)^2=D_\tau,
\]
so this second rational root is also an integer.
Set $f=|f_*|$ and
\[
 r=n+L-f_*.
\]
The strict inequality $L>|f_*|$ gives $r>0$ for both signs of $f_*$.
The first carrier becomes $Z=\overline Z-Mf_*$, so $A_0=0$.
The strong residual is zero because $f^2=D_f$.
All other registers involving the strong root use its square only;
thus the supplied $A_1=A_a=0$ give the remaining weighted-carrier
equations.  The lemma reconstructs a positive source solution.

Conversely, use canonical source witnesses and their positive strong root
$f$.  Choose $T,u$ by positive auxiliary reconstruction, and $r,\gamma$
by the weighted-carrier lemma with $h_*=ikc$.  This multiplier is at least
$k$.  The upper approximation gap satisfies
\[
 0<L-f=\frac{I^2-1}{L+f}<\frac{I}{a}
       =\frac{ic^2}{a}<ikc,
\]
because $c=kY+\eta<k(Y+1)\le ka$ for $X,Y\ge1$.
The transport construction gives
\[
 r>qW-2k>ikc.
\]
The last inequality follows directly from
$q\ge16$ and $ic\ge3$.
Hence $n=r-L+f$ is a strictly positive integer.
Its substitution gives $\overline Z=M(\tau+f)$, and the two source
squares imply $N=0$.  The other two residuals already vanish.
This proves completeness, including positivity of the new coordinate.
\end{proof}

The theorem proves an equivalence and does not decide either side.
The norm is used with its zero-denominator exclusion and both signed-root
branches; an abstract appeal to independent square classes would not
supply these facts.

\subsection{Exact degree, arithmetic cost, and finite certificates}

Write
\[
 q_0=3\beta j,\quad k_0=\eta+\zeta,\quad
 B_{t,0}=(K_0q_0+X)(q_0-F-\alpha)-q_0F.
\]
For the twelve-variable equation,
\[
 \deg M=4,\quad \deg L=7,\quad \deg D_f=14,\quad
 \deg D_\tau=8,\quad \deg Z_0=10.
\]
The cancellation-resolved middle expression in \eqref{eq:rt-N} is
\begin{equation}
 \overline Z^2-M^2(D_f+D_\tau)
 =Z_0^2+2MLZ_0-M^2(1-I^2+D_\tau).\label{eq:rt-centeridentity}
\end{equation}
It has degree $21$; only its middle summand reaches that degree.
The subtracted norm tail $4M^4D_fD_\tau$ has degree $38$.  Thus $N$ has
exact degree $42$, with leading homogeneous form
\[
 \operatorname{in}_{42}(N)
 =4q_0^{16}i^4k_0^8X^2Y^8B_{t,0}^{\,2}.
\]
The other residual degrees are $13$ and $37$, respectively.  Consequently
\begin{equation}
 \operatorname{in}_{84}(P_{12})
 =16q_0^{32}i^8k_0^{16}X^4Y^{16}B_{t,0}^{\,4}.
 \label{eq:rt-leading84}
\end{equation}
In particular the coefficient of
$j^{32}i^8\eta^{16}X^8Y^{16}\alpha^4$ is exactly
$16(3\beta)^{32}>0$.  Choosing this minimum $j$ exponent forces the
$-X\alpha$ term from each copy of $B_{t,0}$, so the other fixed ports
cannot cancel it.  This proves exact degree after the actual coefficient
specialization.

For $P_{13}$, the unique highest residual degree is $15$ and
\[
 \operatorname{in}_{15}(A_0)=-2q_0^8k_0^2XY^2B_{t,0},\qquad
 \operatorname{in}_{30}(P_{13})=4q_0^{16}k_0^4X^2Y^4B_{t,0}^{\,2}.
\]
Its coefficient of $j^{16}\eta^4X^4Y^4\alpha^2$ is
$4(3\beta)^{16}>0$.

The delivered twelve-variable evaluator uses $60$ multiplications,
$27$ additions and $16$ subtractions, totaling $103$ live operations.
The three-equation version uses $98$ operations and residual degrees
$42,13,37$.  The thirteen-variable evaluator uses $94$ operations,
or $87$ for its four-equation system.
The JSON graphs contain no divisions or variable exponents.

The verifier checks literal graph liveness, every fixed-payload digest,
byte-identical circuit regeneration, exact prefix identities,
independent cut-ring norm and finalizer identities, and the nonzero
leading forms above.  These are finite certificates of algebra and data.
The unbounded existence equivalence follows from the written proofs
and the inherited source theorem, not from bounded tests.  No new
end-to-end Lean/Rocq certificate or exported object-language ZFC
derivation is claimed for these literal polynomials.
